% GENERATED FILE. Edit canonical lessons, then run npm run generate:books.
% canonical-source-hash: fnv1a64:e5bdd7aaa7221229
% canonical-lessons: MR-C283-tikit, MR-C283-tikit-dya, MR-C283-haravli, MR-R283-at-the-station

\chapter{Tickets, Times and Trouble}
\label{ch:mr-tickets-times-and-trouble}
\hlchaptermodality{283}

\begin{chapteropening}
\textbf{By the end of this chapter:} \emph{I can buy a ticket (ek tikīṭ dyā), ask when a train leaves, and say that something is lost or missed (haravlī, cuklī).}

Complete MR-R283-at-the-station: buy a ticket, ask the time, and report a problem.
\end{chapteropening}

\section[\texorpdfstring{\mr{तिकीट} (tikīṭ)}{तिकीट (tikīṭ)}]{\mr{तिकीट} (tikīṭ) — a ticket}

\label{lesson:MR-C283-tikit}



\begin{quote}
Before the new one: say \textquotedblleft{}the station is opposite the temple\textquotedblright{} (\emph{sṭeśan mandirācyā samor āhe}), then \textquotedblleft{}a bus\textquotedblright{} and \textquotedblleft{}a train\textquotedblright{}. (\emph{bas}, \emph{gāḍī}.)
\end{quote}

\subsection*{You'll want to know: \mr{तिकीट}}
\textbf{\mr{तिकीट}} — \emph{tikīṭ} — \textquotedblleft{}a ticket\textquotedblright{}.

\begin{quote}
\textbf{\mr{मला} \mr{पुण्याचं} \mr{एक} \mr{तिकीट} \mr{हवं} \mr{आहे}.}
\end{quote}

\begin{quote}
\emph{malā puṇyācaṁ ek tikīṭ havaṁ āhe.} — \textquotedblleft{}I want a ticket to Pune.\textquotedblright{}
\end{quote}

\textbf{\mr{तिकीट}} is neuter, so \textbf{\mr{हवं}}. Two tickets: \textbf{\mr{दोन} \mr{तिकिटं}} (\emph{don tikiṭaṁ}). Marathi buys a ticket with \textbf{\mr{काढणे}}, \textquotedblleft{}to take out\textquotedblright{}: \textbf{\mr{मी} \mr{तिकीट} \mr{काढलं}} (\emph{mī tikīṭ kāḍhlaṁ}), \textquotedblleft{}I bought a ticket\textquotedblright{}.

\begin{cousinweb}
\textbf{\mr{तिकीट}} (\emph{tikīṭ}) is a loan from English \emph{ticket}, reshaped to Marathi sounds: a long \emph{ī} and a retroflex \emph{ṭ}.
\end{cousinweb}

\subsection*{Your turn}
Say these aloud:

\begin{itemize}
\raggedright
  \item \emph{tikīṭ}
  \item \textquotedblleft{}I want a ticket to Pune\textquotedblright{} — \emph{malā puṇyācaṁ ek tikīṭ havaṁ āhe}
  \item \textquotedblleft{}I bought a ticket\textquotedblright{} — \emph{mī tikīṭ kāḍhlaṁ}
\end{itemize}

\begin{tcolorbox}[breakable,colback=teal!4,colframe=teal!35!black,title={Before you move on}]
What does \textbf{\mr{तिकीट}} mean? (\textquotedblleft{}A ticket\textquotedblright{}.) \textquotedblleft{}Two tickets\textquotedblright{}? (\textbf{\mr{दोन} \mr{तिकिटं}}.)
\end{tcolorbox}
\section[\texorpdfstring{\mr{एक} \mr{तिकीट} \mr{द्या} (ek tikīṭ dyā)}{एक तिकीट द्या (ek tikīṭ dyā)}]{\mr{एक} \mr{तिकीट} \mr{द्या} (ek tikīṭ dyā) — at the counter}

\label{lesson:MR-C283-tikit-dya}



\begin{quote}
Ask how far the station is from here. (\emph{ithūn sṭeśan kitī lāmb āhe?}) Then say \textquotedblleft{}I want a ticket to Pune\textquotedblright{}. (\emph{malā puṇyācaṁ ek tikīṭ havaṁ āhe}.) The polite \textquotedblleft{}please give\textquotedblright{} is \textbf{\mr{द्या}}.
\end{quote}

\subsection*{You'll want to know: buying a ticket}
\begin{quote}
\textbf{\mr{पुण्याचं} \mr{एक} \mr{तिकीट} \mr{द्या}.}
\end{quote}

\begin{quote}
\emph{puṇyācaṁ ek tikīṭ dyā.} — \textquotedblleft{}One ticket to Pune, please.\textquotedblright{}
\end{quote}

\textbf{\mr{पुण्याचं}} is \textquotedblleft{}of Pune\textquotedblright{}: the ticket belongs to where it goes. A return: \textbf{\mr{परतीचं} \mr{तिकीट}} (\emph{partīcaṁ tikīṭ}), with the \textbf{\mr{परत}}, \textquotedblleft{}back\textquotedblright{}, of \textbf{\mr{परतणे}}. Then the price: \textbf{\mr{किती} \mr{पैसे} \mr{झाले}?} (\emph{kitī paise jhāle?}), \textquotedblleft{}How much is it?\textquotedblright{}

\subsection*{You'll want to know: when it leaves}
\begin{quote}
\textbf{\mr{पुण्याची} \mr{गाडी} \mr{किती} \mr{वाजता} \mr{निघते}?}
\end{quote}

\begin{quote}
\emph{puṇyācī gāḍī kitī vājtā nighte?} — \textquotedblleft{}What time does the Pune train leave?\textquotedblright{}
\end{quote}

\textbf{\mr{निघते}}, \textquotedblleft{}sets out\textquotedblright{}, is the present of \textbf{\mr{निघणे}}, said of the train. The answers come as times: \textbf{\mr{सात} \mr{वाजता}} (\emph{sāt vājtā}). And the other end: \textbf{\mr{गाडी} \mr{किती} \mr{वाजता} \mr{पोहोचते}?} (\emph{gāḍī kitī vājtā pohocte?}), \textquotedblleft{}What time does it arrive?\textquotedblright{} For the platform, ask with \textbf{\mr{कोणत्या}} (\emph{koṇtyā}, \textquotedblleft{}which\textquotedblright{}): \textbf{\mr{कोणत्या} \mr{फलाटावर}?} (\emph{koṇtyā phalāṭāvar?})

\subsection*{Your turn}
Say these aloud:

\begin{itemize}
\raggedright
  \item \textquotedblleft{}One ticket to Pune, please\textquotedblright{} — \emph{puṇyācaṁ ek tikīṭ dyā}
  \item \textquotedblleft{}A return ticket, please\textquotedblright{} — \emph{partīcaṁ tikīṭ dyā}
  \item \textquotedblleft{}What time does the Pune train leave?\textquotedblright{} — \emph{puṇyācī gāḍī kitī vājtā nighte?}
\end{itemize}

\begin{tcolorbox}[breakable,colback=teal!4,colframe=teal!35!black,title={Before you move on}]
\textquotedblleft{}A ticket to Pune, please\textquotedblright{}? (\textbf{\mr{पुण्याचं} \mr{एक} \mr{तिकीट} \mr{द्या}}.) \textquotedblleft{}A return ticket\textquotedblright{}? (\textbf{\mr{परतीचं} \mr{तिकीट}}.) \textquotedblleft{}What time does it leave?\textquotedblright{} (\textbf{\mr{किती} \mr{वाजता} \mr{निघते}?})
\end{tcolorbox}
\section[\texorpdfstring{\mr{माझी} \mr{पिशवी} \mr{हरवली} (mājhī piśvī haravlī)}{माझी पिशवी हरवली (mājhī piśvī haravlī)}]{\mr{माझी} \mr{पिशवी} \mr{हरवली} (mājhī piśvī haravlī) — when something goes wrong}

\label{lesson:MR-C283-haravli}



\begin{quote}
Say \textquotedblleft{}a bag\textquotedblright{} and \textquotedblleft{}to lose\textquotedblright{}. (\emph{piśvī}, \emph{haravṇe}.) Then \textquotedblleft{}she came\textquotedblright{}. (\emph{tī ālī}.)
\end{quote}

\subsection*{You'll want to know: lost}
\begin{quote}
\textbf{\mr{माझी} \mr{पिशवी} \mr{हरवली}.}
\end{quote}

\begin{quote}
\emph{mājhī piśvī haravlī.} — \textquotedblleft{}My bag is lost.\textquotedblright{}
\end{quote}

The bag is what got lost, so the past agrees with it, as \textbf{\mr{ती} \mr{आली}} agrees with her: \textbf{\mr{पिशवी} \mr{हरवली}}, \textbf{\mr{पुस्तक} \mr{हरवलं}} (\emph{pustak haravlaṁ}). No \textbf{-\mr{ने}} here. Then the report: \textbf{\mr{मी} \mr{पोलिसांकडे} \mr{तक्रार} \mr{केली}} (\emph{mī polisāṅkaḍe takrār kelī}), \textquotedblleft{}I complained to the police\textquotedblright{}, with \textbf{\mr{तक्रार}} feminine.

\subsection*{You'll want to know: missed}
\begin{quote}
\textbf{\mr{माझी} \mr{गाडी} \mr{चुकली}.}
\end{quote}

\begin{quote}
\emph{mājhī gāḍī cuklī.} — \textquotedblleft{}I missed my train.\textquotedblright{}
\end{quote}

Word for word \textquotedblleft{}my train went wrong\textquotedblright{}: with \textbf{\mr{चुकणे}} the train is the subject. The cause comes first: \textbf{\mr{मला} \mr{उशीर} \mr{झाला}, \mr{म्हणून} \mr{माझी} \mr{गाडी} \mr{चुकली}} (\emph{malā uśīr jhālā, mhaṇūn mājhī gāḍī cuklī}), \textquotedblleft{}I was late, so I missed my train\textquotedblright{}.

\subsection*{Your turn}
Say these aloud:

\begin{itemize}
\raggedright
  \item \textquotedblleft{}My bag is lost\textquotedblright{} — \emph{mājhī piśvī haravlī}
  \item \textquotedblleft{}I complained to the police\textquotedblright{} — \emph{mī polisāṅkaḍe takrār kelī}
  \item \textquotedblleft{}I was late, so I missed my train\textquotedblright{} — \emph{malā uśīr jhālā, mhaṇūn mājhī gāḍī cuklī}
\end{itemize}

\begin{tcolorbox}[breakable,colback=teal!4,colframe=teal!35!black,title={Before you move on}]
\textquotedblleft{}My bag is lost\textquotedblright{}? (\textbf{\mr{माझी} \mr{पिशवी} \mr{हरवली}}.) \textquotedblleft{}My book is lost\textquotedblright{}? (\textbf{\mr{माझं} \mr{पुस्तक} \mr{हरवलं}}.) \textquotedblleft{}I missed my train\textquotedblright{}? (\textbf{\mr{माझी} \mr{गाडी} \mr{चुकली}}.)
\end{tcolorbox}
\section[(review)]{Review — at the station}

\label{lesson:MR-R283-at-the-station}



\begin{quote}
Ask the way to the station and how far it is, then give it: straight, then right, opposite the temple. (\emph{sṭeśanlā kasaṁ jāycaṁ? ithūn kitī lāmb āhe?} \emph{saraḷ jā, mag ujvīkaḍe vaḷā. sṭeśan mandirācyā samor āhe.}) Then say \textbf{\mr{ऋतू}} and \textbf{\mr{ज्ञान}}, and name the letters \textbf{\mr{ऋ}} and \textbf{\mr{ञ}}.
\end{quote}

\subsection*{Your turn}
Say these aloud:

\begin{itemize}
\raggedright
  \item One ticket to Pune, please. A return ticket.
  \item What time does the train leave? Which platform?
  \item My bag is lost. I complained to the police.
  \item I was late, so I missed my train.
  \item what went wrong on a trip of yours, in two sentences
\end{itemize}

\begin{tcolorbox}[breakable,colback=teal!4,colframe=teal!35!black,title={Before you move on}]
A ticket? (\textbf{\mr{तिकीट}}.) \textquotedblleft{}When does it leave?\textquotedblright{} (\textbf{\mr{किती} \mr{वाजता} \mr{निघते}?}) \textquotedblleft{}Lost\textquotedblright{}? (\textbf{\mr{हरवली}}.) \textquotedblleft{}Missed\textquotedblright{}? (\textbf{\mr{चुकली}}.) Which letters hide in \textbf{\mr{ज्ञान}} and start \textbf{\mr{ऋतू}}? (\textbf{\mr{ञ}}, \textbf{\mr{ऋ}}.)
\end{tcolorbox}
